% Источники: raw/IMG_7629.jpeg -- raw/IMG_7635.jpeg.
% Компиляция: tectonic 07.tex
\ifdefined\lessonpreambleloaded
  \expandafter\endinput
\fi
\let\lessonpreambleloaded\empty
\documentclass[a4paper,11pt]{article}
\usepackage{fontspec}
\setmainfont{Times New Roman}
\usepackage[russian]{babel}
\usepackage{amsmath,amssymb}
\usepackage[margin=22mm]{geometry}
\usepackage{indentfirst}
\usepackage{needspace}
\usepackage{tikz}
\usetikzlibrary{patterns}
\usepackage[hidelinks]{hyperref}
\usepackage{fancyhdr}
\pagestyle{fancy}
\fancyhf{}
\renewcommand{\headrulewidth}{0pt}
\fancyfoot[C]{\footnotesize\url{https://github.com/bmstu-software-engineering/cs-s5-tv-26y}}
\fancyfoot[R]{\thepage}
\setlength{\emergencystretch}{2em}
\newcommand{\sourcepage}[1]{\par\Needspace{6\baselineskip}\medskip\noindent{\footnotesize\itshape Источник: \detokenize{#1}.}\par\nopagebreak\smallskip}
\newcommand{\editorial}[1]{\par\smallskip\noindent{\small\itshape Примечание к переносу. #1}\par\smallskip}
\newcommand{\lessontitle}[3]{\begin{center}{\large\bfseries Теория вероятностей. #1}\\[4pt]{\large #2}\\[4pt]#3\end{center}}
\newcommand{\parttitle}[1]{\par\Needspace{4\baselineskip}\medskip\noindent\textbf{#1}\par\nopagebreak\smallskip}
\tikzset{element/.style={draw,rectangle,minimum width=7mm,minimum height=5mm,fill=white},every picture/.style={line width=.5pt}}

\begin{document}
\lessontitle{Семинар №\,7}{Комбинаторика. Сложение и умножение вероятностей}{17.10.26}

\sourcepage{IMG_7629.jpeg}
\section*{Полиномиальный коэффициент}
\[
 C(m_1,m_2,\ldots,m_k)=\frac{n!}{m_1!m_2!\cdots m_k!}.
\]
Здесь $n$ --- число элементов, $k$ --- число классов (групп),
$m_i$ --- число элементов в $i$-й группе; $m_1+\cdots+m_k=n$.
Порядок групп важен, порядок элементов внутри группы не важен.
Биномиальный коэффициент --- обычный коэффициент сочетаний.

\textbf{Пример.} Десять студентов распределяются по трём помеченным группам:
в первую --- $5$, во вторую --- $3$, в третью --- $2$:
\[
 C(5,3,2)=\frac{10!}{5!3!2!}=2520.
\]
Внизу страницы дополнительно записано:
\[
 P_3=3!=6,\qquad \frac{C(5,3,2)}{6}=420.
\]
\editorial{Условие, к которому относится деление на $6$, не приведено.
При помеченных группах ответом остаётся $2520$; деление на $6$ здесь не обосновывается.}

\sourcepage{IMG_7630.jpeg}
\section*{Телефонные каналы}
На узел связи поступили шесть телефонных вызовов; они распределяются
по четырём каналам независимо друг от друга. Рассматривается распределение:
на первый канал --- $3$ вызова, на второй --- $2$, на третий --- $1$,
на четвёртый --- $0$.

Исход $(x_1,\ldots,x_6)$, где $x_i\in\{1,2,3,4\}$ --- номер канала
$i$-го вызова. Это размещения с повторениями:
\[
 \Omega=\{(x_1,\ldots,x_6):x_i\in\{1,2,3,4\}\},\qquad
 N=\widetilde A_4^6=4^6.
\]
Событие $A$ задаётся условиями
\[
 |\{i:x_i=1\}|=3,\quad |\{i:x_i=2\}|=2,\quad
 |\{i:x_i=3\}|=1,\quad |\{i:x_i=4\}|=0.
\]
\[
 N_A=C(3,2,1,0)=\frac{6!}{3!2!1!0!},\qquad P(A)=\frac{N_A}{N}.
\]
\editorial{В исходнике итоговая дробь не вычислена.
Равновозможность каналов используется в записанном подсчёте $N=4^6$.}

\parttitle{Домашнее задание}
Номер 2.28. Краткие записи событий:
$A$ --- разные вагоны; $B$ --- один вагон пуст;
$C$ --- в I вагоне $5$, во II --- $2$, в III --- $1$, в IV --- $1$,
остальные пусты.
\editorial{Полное условие домашней задачи на странице не приведено.}

\sourcepage{IMG_7631.jpeg}
\section*{Сложение и умножение вероятностей}
\begin{enumerate}
 \item $0\le P(A)\le1$.
 \item $P(\overline A)=1-P(A)$.
 \item Если $A$ и $B$ несовместны, то
 \[
 P(A\cup B)=P(A)+P(B).
 \]
 Несовместные события не могут произойти одновременно:
 $A\cap B=\varnothing$. Для произвольных событий
 \[
 P(A\cup B)=P(A)+P(B)-P(A\cap B).
 \]
 \item Если $A$ и $B$ независимы, то
 \[
 P(A\cap B)=P(A)P(B).
 \]
 Независимость означает, что наступление одного события не меняет
 вероятность другого. В общем случае при $P(A)>0$
 \[
 P(A\cap B)=P(A)P(B\mid A).
 \]
\end{enumerate}
В исходнике используются обозначения $A+B=A\cup B$ и $AB=A\cap B$.
\begin{center}
\begin{tikzpicture}[scale=.75]
 \fill[gray!20] (0,0) circle (1);\fill[gray!20] (1.3,0) circle (1);
 \draw (0,0) circle (1) (1.3,0) circle (1);
 \node at (-.45,0){$A$};\node at (1.75,0){$B$};\node at (.65,-1.35){$A\cup B$};
 \begin{scope}[xshift=5.5cm]
 \begin{scope}\clip (0,0) circle (1);\fill[gray!40] (1.3,0) circle (1);\end{scope}
 \draw (0,0) circle (1) (1.3,0) circle (1);
 \node at (-.45,0){$A$};\node at (1.75,0){$B$};\node at (.65,-1.35){$A\cap B$};
 \end{scope}
\end{tikzpicture}
\end{center}

\sourcepage{IMG_7632.jpeg}
\section*{Обнаружение воздушной цели}
Две станции независимо обнаруживают цель. Вероятности обнаружения:
$0{,}7$ для первой и $0{,}8$ для второй. Найти вероятность обнаружения
ровно одной станцией; затем --- хотя бы одной.

$A$ --- цель обнаружена первой станцией, $B$ --- второй:
\[
 P(A)=0{,}7,\quad P(\overline A)=0{,}3,\quad
 P(B)=0{,}8,\quad P(\overline B)=0{,}2.
\]
\parttitle{1. Ровно одна станция}
\[
 C=(A\cap\overline B)\cup(\overline A\cap B).
\]
Слагаемые события несовместны. Используя независимость, получаем
\[
 P(C)=P(A)P(\overline B)+P(\overline A)P(B).
\]
\parttitle{2. Хотя бы одна станция}
В исходнике записаны три способа описать событие $D$:
\[
 D=(A\cap\overline B)\cup(\overline A\cap B)\cup(A\cap B),
\]
\[
 P(D)=P(A\cup B)=P(A)+P(B)-P(A)P(B),
\]
\[
 \overline D=\overline A\cap\overline B,\qquad
 P(\overline D)=0{,}06,\qquad P(D)=1-0{,}06=0{,}94.
\]

\sourcepage{IMG_7633.jpeg}
\section*{Задачи с цепями}
Рассматриваются последовательное и параллельное соединения.
$A$ --- система работает; $\overline A$ --- отказ системы;
$A_i$ --- работает $i$-й элемент, $\overline A_i$ --- его отказ.
Требуется описать $A$, $\overline A$ и найти $P(A)$.

\parttitle{Последовательное соединение}
\begin{center}
\begin{tikzpicture}
 \draw (0,0)--(5,0);
 \node[element] at (1,0){1};\node[element] at (2.5,0){2};\node[element] at (4,0){3};
 \fill (0,0) circle (1.3pt);\fill (5,0) circle (1.3pt);
\end{tikzpicture}
\end{center}
Система работает, если работают все элементы:
\[
 A=A_1\cap A_2\cap A_3,\qquad
 \overline A=\overline A_1\cup\overline A_2\cup\overline A_3.
\]
При независимых отказах
\[
 P(A)=P(A_1)P(A_2)P(A_3),\qquad
 P(\overline A)=1-P(A_1)P(A_2)P(A_3).
\]
\parttitle{Параллельное соединение}
\begin{center}
\begin{tikzpicture}
 \draw (0,0)--(1,0) (1,-.8)--(1,.8) (3,-.8)--(3,.8) (3,0)--(4,0);
 \foreach \y/\n in {.8/1,0/2,-.8/3}{\draw (1,\y)--(3,\y);\node[element] at (2,\y){\n};}
 \fill (0,0) circle (1.3pt);\fill (4,0) circle (1.3pt);
\end{tikzpicture}
\end{center}
\[
 A=A_1\cup A_2\cup A_3,\qquad
 \overline A=\overline A_1\cap\overline A_2\cap\overline A_3.
\]
\[
 P(A)=1-P(\overline A)=1-P(\overline A_1)P(\overline A_2)P(\overline A_3).
\]

\sourcepage{IMG_7634.jpeg}
\section*{Система из четырёх элементов}
Элементы выходят из строя независимо. Даны вероятности работы:
\[
 p_1=0{,}6,\quad p_2=0{,}8,\quad p_3=0{,}9,\quad p_4=0{,}7.
\]
Вероятности отказа: $q_1=0{,}4$, $q_2=0{,}2$, $q_3=0{,}1$, $q_4=0{,}3$.
На схеме верхняя ветвь содержит последовательно элементы 1 и 2;
ещё две параллельные ветви содержат элементы 3 и 4.
\begin{center}
\begin{tikzpicture}
 \draw (0,0)--(1,0) (1,-.9)--(1,.9) (4,-.9)--(4,.9) (4,0)--(5,0);
 \foreach \y in {-.9,0,.9}{\draw (1,\y)--(4,\y);}
 \node[element] at (2,.9){1};\node[element] at (3,.9){2};
 \node[element] at (2.5,0){3};\node[element] at (2.5,-.9){4};
\end{tikzpicture}
\end{center}
\[
 A=(A_1\cap A_2)\cup(A_3\cup A_4),
 \qquad \overline A=(\overline A_1\cup\overline A_2)
 \cap\overline A_3\cap\overline A_4.
\]
\[
 P(\overline A)=(1-p_1p_2)q_3q_4=0{,}0156,
 \qquad P(A)=1-P(\overline A)=0{,}9844.
\]

\sourcepage{IMG_7635.jpeg}
\section*{Дополнительные схемы}
\parttitle{Три параллельных элемента и один последовательный}
\begin{center}
\begin{tikzpicture}
 \draw (0,0)--(1,0) (1,-.8)--(1,.8) (3,-.8)--(3,.8) (3,0)--(5,0);
 \foreach \y/\n in {.8/1,0/2,-.8/3}{\draw (1,\y)--(3,\y);\node[element] at (2,\y){\n};}
 \node[element] at (4,0){4};
\end{tikzpicture}
\end{center}
\[
 A=(A_1\cup A_2\cup A_3)\cap A_4,\qquad
 \overline A=(\overline A_1\cap\overline A_2\cap\overline A_3)\cup\overline A_4,
\]
\[
 P(A)=(1-q_1q_2q_3)p_4.
\]
\parttitle{Параллельные 1 и 2 последовательно с 3; нижняя ветвь --- 4}
\begin{center}
\begin{tikzpicture}
 \draw (0,0)--(1,0) (1,-.8)--(1,1.2) (3,.4)--(3,1.2)
 (1,1.2)--(3,1.2) (1,.4)--(3,.4) (3,.8)--(4,.8)
 (4,.8)--(4,-.8) (1,-.8)--(4,-.8) (4,0)--(5,0);
 \node[element] at (2,1.2){1};\node[element] at (2,.4){2};
 \node[element] at (3.5,.8){3};\node[element] at (2.5,-.8){4};
\end{tikzpicture}
\end{center}
\[
 A=((A_1\cup A_2)\cap A_3)\cup A_4,
 \qquad\overline A=((\overline A_1\cap\overline A_2)\cup\overline A_3)\cap\overline A_4,
\]
\[
 P(\overline A)=\bigl(1-(1-q_1q_2)p_3\bigr)q_4.
\]
\editorial{В последней формуле рукописи у вероятности отказа четвёртого элемента
пропущено дополнение; здесь использовано $q_4=P(\overline A_4)$.}
Последняя схема помечена как домашнее задание.
\end{document}
